\section{Mathematical Formulation and Implementation}
\label{sec:formulation}

\subsection{Decision Variables}

Let $x_{r,t} \geq 0$ denote the compute load (kWh) assigned to region $r \in \mathcal{R}$ at hour $t \in \mathcal{T}$.
This single set of variables simultaneously encodes both temporal scheduling (varying $t$ for a fixed $r$) and spatial routing (varying $r$ for a fixed $t$), enabling the model to find a globally optimal joint allocation without imposing a sequential two-phase structure.

\subsection{Objective Function}

The model minimizes total carbon emissions, adjusted to incentivize consumption during high-renewable periods:

\begin{equation}
  \min_{x} \quad \sum_{r \in \mathcal{R}} \sum_{t \in \mathcal{T}}
    x_{r,t} \cdot \mathrm{CI}_r(t) \cdot \bigl(1 - \alpha \cdot \mathrm{RF}_r(t)\bigr)
  \label{eq:objective}
\end{equation}

\textbf{Interpretation of terms.}
The factor $\mathrm{CI}_r(t)$ is the baseline carbon cost per unit of compute~\cite{radovanovic2023}.
The weighting term $(1 - \alpha \cdot \mathrm{RF}_r(t))$ provides an additional discount when the renewable fraction is high: a low-CI hour driven by wind and solar (high RF) receives a stronger incentive than a low-CI hour driven by nuclear (low RF).
This distinction is important because absorbing variable renewable surplus reduces the need for curtailment and supports grid decarbonization beyond what CI alone captures~\cite{liu2012, riepin2024spatiotemporal}.

\textbf{Tuning parameter $\alpha$.}
$\alpha \in [0, 1]$ controls the strength of the RF incentive.
At $\alpha = 0$, the model reduces to pure CI minimization.
Sensitivity analysis over $\alpha$ is reported in Section~\ref{sec:results}.

\subsection{Constraints}

\begin{align}
  &\sum_{r \in \mathcal{R}} \sum_{t \in \mathcal{T}} x_{r,t} = D
  \tag{C1} \label{eq:c1} \\[4pt]
  &x_{r,t} \leq C_{\max,r} \quad \forall\, r \in \mathcal{R},\; t \in \mathcal{T}
  \tag{C2} \label{eq:c2} \\[4pt]
  &x_{\mathrm{flex},r,t} \leq \mathrm{VCC}_r(t) \quad \forall\, r \in \mathcal{R},\; t \in \mathcal{T}
  \tag{C3} \label{eq:c3} \\[4pt]
  &\sum_{t \in \mathcal{T}} x_{r,t} \cdot P_r(t) \leq B_r \quad \forall\, r \in \mathcal{R}
  \tag{C4} \label{eq:c4} \\[4pt]
  &x_{r,t} \geq 0 \quad \forall\, r \in \mathcal{R},\; t \in \mathcal{T}
  \tag{C5} \label{eq:c5}
\end{align}

\textbf{C1 — Demand satisfaction.}
Total compute demand $D$ must be fulfilled exactly over the planning horizon.
This is a hard constraint; no load shedding is permitted~\cite{liu2012}.

\textbf{C2 — Regional capacity.}
At each hour, region $r$ cannot accept more than $C_{\max,r}$ kWh of compute load, reflecting physical server capacity and power supply limits~\cite{liu2012}.

\textbf{C3 — Virtual Capacity Curve.}
The flexible component of load in region $r$ at hour $t$ is bounded by $\mathrm{VCC}_r(t)$, a CI-indexed ceiling derived from the supervisor's grid-interactive datacenter framework~\cite{colangelo2025}.
The VCC encodes the operator's commitment to reduce load when the grid is under carbon stress.

\textbf{C4 — Per-region electricity budget.}
The electricity cost in each region $r$ must remain within budget $B_r$, where $P_r(t)$ is the spot electricity price at hour $t$~\cite{hao2024joint}.
Since PJM, NYISO, Nord Pool, and USEP prices are denominated in different currencies, this constraint is applied as a normalized per-region cap rather than a cross-region cost comparison.

\textbf{C5 — Non-negativity.}
Load assignments are non-negative by construction.

\subsubsection*{Component-Level References}

\begin{table}[h]
\centering
\begin{tabular}{lp{7cm}l}
\toprule
\textbf{Component} & \textbf{Description} & \textbf{Reference} \\
\midrule
Objective \eqref{eq:objective}
  & CI-weighted load as carbon cost signal
  & \citet{radovanovic2023} \\
RF weighting $(1-\alpha\cdot\mathrm{RF})$
  & Novel extension distinguishing renewable source quality
  & This thesis \\
Constraints \eqref{eq:c1}--\eqref{eq:c2}
  & Demand completion and regional capacity bounds
  & \citet{liu2012} \\
Constraint \eqref{eq:c3}
  & VCC-indexed flexible load ceiling
  & \citet{colangelo2025} \\
Constraint \eqref{eq:c4}
  & Per-region electricity cost budget
  & \citet{hao2024joint} \\
Joint spatio-temporal LP structure
  & Cross-region temporal-spatial allocation in one model
  & \citet{riepin2024spatiotemporal, attenni2024shifting} \\
\bottomrule
\end{tabular}
\caption{Component-level references for the mathematical formulation.}
\label{tab:modelrefs}
\end{table}

\subsection{Implementation}

\subsubsection*{Solver}
The model is implemented in Python using the \texttt{gurobipy} interface to Gurobi~10.
With $|\mathcal{R}| = 5$ regions and $|\mathcal{T}| = 17{,}544$ hours, the LP has 87,720 continuous variables and approximately 105,264 constraints.
At this scale, Gurobi's simplex solver typically reaches optimality in under a minute for a single full-horizon solve.
\todo{Report actual solve time after implementation.}

\subsubsection*{Rolling-window backtesting}
To simulate real-world scheduling, the LP is solved in a rolling-window fashion with a decision horizon of $W$ hours and a look-ahead of $L$ hours.
\todo{Specify W and L; report sensitivity to window size.}

\subsubsection*{Data pipeline}
CI and RF inputs are loaded from pre-computed Parquet files (ElectricityMaps, 2024--2025).
Electricity price data is \todo{sourced from: specify}.
All inputs are aligned to hourly UTC timestamps before passing to the solver.

\subsection{Heuristic Fallback}

If rolling-window LP solve times prove prohibitive, a greedy carbon-aware heuristic is used as a fallback and as an additional baseline:

\begin{algorithm}
\caption{Greedy Carbon-Aware Shifting}
\begin{algorithmic}[1]
\REQUIRE Remaining demand $d$, window $[t_1, t_2]$, regions $\mathcal{R}$, CI, RF, capacities
\STATE Compute effective cost $\hat{c}_{r,t} = \mathrm{CI}_r(t) \cdot (1 - \alpha \cdot \mathrm{RF}_r(t))$ for all $(r, t)$
\STATE Sort all $(r, t)$ pairs by $\hat{c}_{r,t}$ ascending
\FOR{each $(r, t)$ in sorted order}
  \STATE Assign $x_{r,t} = \min(d,\; C_{\max,r},\; \mathrm{VCC}_r(t))$
  \STATE $d \leftarrow d - x_{r,t}$
  \IF{$d = 0$} \STATE \textbf{break} \ENDIF
\ENDFOR
\end{algorithmic}
\end{algorithm}

The heuristic is polynomial-time and provides a near-optimal lower bound on carbon cost, useful for benchmarking LP solution quality.
